\dresumo{
Organizações que acumularam sistemas ao longo do tempo frequentemente dependem de informações dispersas em várias fontes, o que exige consulta e consolidação manuais e aumenta o risco de retrabalho e inconsistências. Este trabalho analisa o InLight, aplicativo desenvolvido para unificar as informações do BackOffice de compra de energia do Ambiente de Contratação Regulada (ACR) da Light, e propõe sua implantação e avaliação, considerando os benefícios e os impactos da unificação de sistemas corporativos. A pesquisa, aplicada, exploratória e descritiva, adota o estudo de caso único e combina uma fundamentação teórica sobre sistemas de informação, integração, processos de negócio e aceitação de tecnologia com uma revisão da literatura realizada na base Scopus, que resultou em 11 estudos sintetizados. A partir desse referencial, o trabalho diagnostica o cenário atual da área, no qual as informações estão distribuídas entre os sistemas da CCEE, a plataforma Thunders, um sistema de lançamento de notas e planilhas de controle. Descreve os módulos, a arquitetura e as tecnologias do InLight propõe um plano de implantação baseado nas fases do ciclo de vida de sistemas corporativos e em fatores críticos de sucesso e define indicadores de processo, um questionário de percepção dos usuários e um roteiro de entrevista com gestores para avaliar seus resultados. Como o aplicativo ainda não foi implantado, os benefícios esperados são apresentados como hipóteses de trabalho. Conclui-se que os ganhos da unificação dependem de a solução responder a problemas concretos do processo, de o processo ser revisto junto com a tecnologia e de os usuários perceberem valor no novo modo de trabalho.
}
{unificação de sistemas; sistemas corporativos; implantação de sistemas; setor elétrico; ambiente de contratação regulada.}
